\documentclass[../dayan_kernel_manual.tex]{subfiles}

\begin{document}

\section{算子编程模型}\label{sec:kernel-010}

dycc 源文件可以同时包含 Host 和 Device 代码。每个算子源文件显式包含
\code{dayan_kernel.h}；Device API 头文件放在 \code{__DAYAN_DEVICE__} 分支内。

\subsection{源码切分}\label{sec:kernel-022}

\begin{dayantable}{源码属性}{tab:kernel-010}{X[20,l] X[26,l] X[26,l] X[28,l]}
  属性 & Host & Device & 规则 \\
  \code{__kernel__} & 生成同名 wrapper & 保留函数 & 仅用于 kernel 函数。 \\
  \code{__host__} & 保留 & 删除 & 用于 Host 函数或全局变量。 \\
  \code{__device__} & 删除 & 保留 & 用于 Device 函数或全局变量。 \\
  无属性函数 & 保留 & 保留 & 声明和函数体必须能在两侧编译。 \\
  无属性全局变量 & 报错 & 报错 & 全局变量必须指定所属侧。 \\
\end{dayantable}

\begin{dayancode}
#include <dayan_kernel.h>

#if defined(__DAYAN_DEVICE__)
#include <dayan_sys.h>
#include <stdio.h>
#endif
#include <stdint.h>

__host__ uint64_t host_bias = 1;
__device__ uint64_t device_calls = 0;

uint64_t clamp_count(uint64_t n)
{
    return n > 1024 ? 1024 : n;
}

__device__ uint64_t next_call(void)
{
    return ++device_calls;
}

__kernel__ void demo_kernel(uint64_t count)
{
    printf("count=%lu call=%lu\n", (unsigned long)clamp_count(count),
           (unsigned long)next_call());
    fflush(stdout);
    dayan_yield();
}

DAYAN_KERNEL_DESC(demo_kernel, 1, 1024, 32768, 0, 0, 0, 0);
\end{dayancode}

\code{main} 和只使用 Host 库的辅助函数应标为 \code{__host__}。同一条逗号声明中的
全局变量不能分属两侧。需要切除的函数或变量定义不能来自宏展开，因为 dycc 必须按
源码范围改写声明。

\subsection{kernel 接口}

\code{__kernel__} 函数应满足：

\begin{itemize}
  \item 返回 \code{void}，具有外部链接，并在当前翻译单元中恰有一个定义；
  \item 不是函数模板或可变参数函数；
  \item 参数不超过 8 个，类型为整数、枚举、指针、引用、\code{float} 或
        \code{double}。
\end{itemize}

聚合体不能按值传入；把数据放入 Runtime Buffer。Device 侧的指针或引用在 Host
wrapper 中改写为 \code{npu_buffer_t *}，应用不会取得或传入 Device 物理地址。

Host 侧的同名 wrapper 返回状态，并在参数末尾追加启动属性。例如：

\begin{dayancode}
/* Device entry: void demo_kernel(uint64_t *, uint64_t); */
int demo_kernel(npu_buffer_t *output, uint64_t count,
                const npu_kernel_launch_attrs_t *attrs);
\end{dayancode}

整数、枚举和地址直接占用一个 64-bit 参数槽；\code{float} 和 \code{double} 传递
位模式。wrapper 使用线程安全的 static descriptor 保存嵌入对象和 C++ mangled entry，
Device module 则按 Context 缓存。首次调用通过
\code{npu_kernel_descriptor_create} 完成描述符创建，后续调用走原子快路径。

每个 kernel 还必须定义
\code{DAYAN_KERNEL_DESC(kernel, child\_threads, atomic\_bytes, uram\_bytes,
lsu\_tcsr, mat\_tcsr, vec\_tcsr, macro\_tcsr)}。
dycc 在 Device object 中保留外部 C 符号 \code{<kernel>\_\_desc}，并把相同常量写入
Host Package。资源值表示一个实例允许同时存活的子 Thread 上限、由
\code{shared/gp} 寻址的 atomic-cache allocation、本 Core 私有 URAM，以及四类 TCSR
池大小；URAM 必须是 32 KiB 的整数倍。子 Thread 上限为 15；root 也占一个 proc slot，
Runtime 保证同一 Hart 上所有 resident invocation 的 \code{1 + child\_threads} 总和不
超过 16，并静态分配 TCSR 物理范围。

允许最多 1--4 个 Task 共驻同一 Hart 时使用
\code{DAYAN_KERNEL_TASK_DESC(kernel, tasks\_per\_core, child\_threads,
atomic\_bytes, uram\_bytes, lsu\_tcsr, mat\_tcsr, vec\_tcsr, macro\_tcsr)}。
TCSR 已由 Context 保存和映射；如果 single-Core kernel 还依赖其它未保存的 Hart 状态，
则必须使用 \code{DAYAN_KERNEL_LOCKED_TASK_DESC} 的同类资源参数。一个带锁 logical Task
只能进入空 Hart；在其整个 Proc family 结束前，
其它 logical Task（包括同一 grid 的下一 Task）不能进入该 Hart，而该 Task 自己的 child
Thread 仍可在声明的 slot 内运行。锁定属性独立于 Thread 数量，不应通过虚报
\code{child\_threads=15} 代替。

\subsection{算子分片与 Proc 共驻}

激活、Element-wise 以及其它分片间没有依赖的算子，默认应采用无子 Thread
的顺序 Proc 模型。Host 先把输出坐标集合分成足够多的逻辑分片，每个 Task
坐标对应一个根 Proc。Proc 根据 \code{dayan\_task\_id/dim} 计算自己的分片范围，
然后在本地循环中顺序执行
\code{Load tile -> compute -> Store tile}。每个阶段只需在相应执行单元上完成后
调用 \code{dayan\_wait\_*}，不需要为 Load、计算和 Store 分别创建 Thread。

这类 kernel 通常使用
\code{DAYAN\_KERNEL\_TASK\_DESC(kernel, 4, 0, \ldots)}，把每 Hart 最大共驻数声明为
4。描述符中的 URAM 和 LSU/MAT/VEC/MACRO TCSR 是单个 Proc 的最坏资源；
Server 只在 proc slot、URAM 和四类 TCSR 都能静态满足时才实际共驻四个 Proc。
这是入场上限而不是固定绑核承诺；当某类资源不足时，Scheduler 自然降低实际
并发数。

当一个 Proc 在 LSU、MAT、VEC 或 MACRO 完成点等待时，协作式 Scheduler 会运行
同 Hart 上其它就绪 Proc。因此多个顺序 Proc 的 Load、计算和 Store 可在不同
执行单元上交错推进，用较小的软件复杂度提高利用率。每个 Proc 有独立 Context，
所以仍应优先使用逻辑 TCSR 组 0。

子 Thread 只用于数据复用或同步需求明确超过上述模型的算子。典型场景是四核矩阵乘：
\code{kFourCore} 的四个 Proc 通过 coherent shared 状态协作，在 Core 间广播 A/B tile，
并在每个 Core 内创建 loader、remote-copy、MAT 或 store Thread 组成生产者/消费者
流水。只有这类收益经验证的特殊算子才应承担 Thread、resource 和 join 的额外开销。

\subsection{启动 kernel}\label{sec:kernel-023}

BL1 与 API ELF 由独立的 \code{dayan_server} 初始化。应用只连接 Service 并创建逻辑
对象：

\begin{dayancode}
npu_runtime_t *runtime;
npu_context_t *context;
npu_stream_t *stream;
npu_buffer_t *output;

npu_runtime_open(&runtime, nullptr);
npu_context_create(runtime, nullptr, &context);
npu_stream_create(context, nullptr, &stream);
npu_buffer_desc_t buffer_desc;
npu_buffer_desc_init(&buffer_desc, 4096);
npu_buffer_create(context, &buffer_desc, &output);

npu_kernel_launch_attrs_t attrs;
npu_kernel_launch_attrs_init(&attrs, context, stream, "demo-task");
attrs.config.task_dim = {8, 1, 1};
attrs.config.task_mode = NPUTaskMode::kSingleCore;

int status = demo_kernel(output, 64, &attrs);
\end{dayancode}

\begin{dayantable}{启动属性}{tab:kernel-011}{X[32,l] X[68,l]}
  字段 & 含义 \\
  \code{version}、\code{reserved} & 由初始化函数设置；reserved 保持 0。 \\
  \code{context}、\code{stream} & 启动所属资源域和有序命令队列；二者必须属于同一 Context。 \\
  \code{task_name} & 本次 launch 的 root proc 名称；必须非空且最多 15 bytes。 \\
  \code{config.task_dim} & x、y、z 三维 logical task 数量，三个分量都必须非零。 \\
  \code{config.task_mode} & \code{kSingleCore} 或 \code{kFourCore}；不含物理 Cluster/Core 编号。 \\
  \code{flags} & 0 为同步；\code{NPU_KERNEL_LAUNCH_ASYNC} 为异步。 \\
  \code{timeout_ms} & 同步等待超时；0 使用 Runtime 默认值。 \\
  \code{cookie} & 原样记录到表示整个 grid 的 Event。 \\
  \code{event_out} & 异步模式的 Event 输出；异步启动时不能为空。 \\
\end{dayantable}

同步 wrapper 等待并销毁 Event，返回整个 grid 的聚合状态。异步 wrapper 返回后，
调用者使用 \code{npu_event_query}、\code{npu_event_wait} 和
\code{npu_event_destroy} 管理 \code{event_out}。

\subsection{logical grid 与调度模式}

\code{NPUTaskDim3} 与 Triton launch grid 类似，只描述 x、y、z 三个方向上的 logical
task 数量。\code{kSingleCore} 以单 Core task 为调度单元，Service 展平并调度
\code{x*y*z} 个 task。\code{kFourCore} 以一个四 Core Cluster 为调度单元，要求
\code{task_dim.x == 4}；每个 $(y,z)$ group 原子占用同一 Cluster 的 Core 0--3，
x 坐标分别为 0--3。

Runtime、Service 和 BL1 都校验 mode 与维度。kernel 用
\code{dayan_task_id}、\code{dayan_task_dim} 和 \code{dayan_task_mode} 读取逻辑坐标，
不会看到物理 Cluster/Core。

\subsection{Runtime 对象}

Context 是一个客户端在逻辑 Device 上的资源域，不等同于 OS 线程、进程或物理
Cluster。Buffer、Stream、Module、Kernel 和 Event 都属于一个 Context；跨 Context
组合对象会被拒绝。销毁 Context 会取消排队任务、等待已经授权的数据 token 收尾，并
统一释放子对象和 Device allocation。

Stream 是有序命令队列。同一 Stream 的传输、launch、wait 和 Event 按提交顺序推进；
不同 Stream 可以并发，但不保证同时占用硬件。\code{npu_stream_wait_event} 建立显式的
跨 Stream 依赖；Host 阻塞由 \code{npu_event_wait} 或
\code{npu_stream_synchronize} 表达。

Module 是已加载的 Kernel Package，Kernel 是其中一个可启动入口。Package 元信息包含
目标 ISA、参数布局、支持的 task mode、每 Core 最大 task 数、resident memory、
子 Thread/atomic cache/URAM/TCSR 上限和同步资源声明。不可变代码和只读数据允许跨 Context
共享；可写数据、BSS 和 Context 相关
重定位结果保持隔离。同步资源只描述类型、数量和初始参数，地址由 Service/HAL 分配。

Buffer 是 Runtime 管理的逻辑 Device 内存对象。应用只使用 Buffer handle 和 offset；
Service 根据可见性选择调度目标，HAL 决定具体 region、地址和传输 provider。

\subsection{执行模型}

每个 hart 运行独立的协作式 scheduler。kernel 对应一个 proc，只有在调用 yield、
等待执行单元、等待 resource、执行阻塞 RPC 或退出时才切回 scheduler。Host mailbox
使用轮询状态，不会抢占正在执行的 proc。

\code{dayan_yield} 只把当前 proc 放回本 hart 的 ready queue。系统没有优先级、时间片
或跨 hart 迁移；长时间不进入等待点的 kernel 会阻塞同一 hart 上的其它 proc。

LSU、MAT、VEC 和 MACRO 可以异步执行。\code{dayan_wait_*} 等待调用前已发射到对应
单元的工作；\code{dayan_resource_wait_*} 与
\code{dayan_shared_resource_wait_*} 则非阻塞地把完成事件绑定到一次资源
produce/consume。硬件 Produce 队列同时保证全局 32 项、每单元 2 项；软件不能超过这一
无阻塞容量，异步事件引用的资源在完成前必须保持有效。

\subsection{内存模型}

Device RAM 中的普通数据、栈和私有堆属于单个 hart 的非一致性访问域。跨 hart 数据
放在 Host coherent window 中；当前 proc 的 allocation 基址由 \code{shared} 读取：

\begin{dayancode}
#include <dayan_sys.h>

struct layout *data = (struct layout *)shared;
\end{dayancode}

\code{shared == NULL} 表示没有 coherent allocation。BL1 在 proc 中保存 allocation
长度，DAPI 会据此检查 shared resource 和 barrier 的精确边界；Host 和算子仍必须使用
同一布局。\code{shared} 绑定 \code{gp}，只读其值，不要给该变量赋值。

URAM 是每 Core 私有空间。kernel 用 \code{dayan\_uram\_base/size} 读取本实例范围；
默认容量为 1.5 MiB、Core 地址窗口 stride 为 2 MiB，分配粒度为 32 KiB。子 Thread
继承根 Proc 的 shared 与 URAM 范围，不能跨 Core 访问其它实例的 URAM；TCSR 必须在
创建 Thread 时显式声明，并通过独立 Context 从逻辑组 0 开始映射。

每 hart 私有堆供 \code{malloc} 使用。初始 region 不足时，Dayan OS 通过 Host RPC
扩展；私有堆地址不能跨 hart 释放，也不能用于跨核共享。

\subsection{resource 模型}

两类 resource 都遵循 \code{FREE -> BUSY -> READY -> FREE}。\code{alloc} 设置
producer/consumer 数量；最后一个 producer 发布 READY，最后一个 consumer 释放为
FREE。

\begin{dayantable}{resource 类型}{tab:kernel-012}{X[29,l] X[31,l] X[40,l]}
  类型 & 存储位置 & 使用范围 \\
  \code{dayan_resource_t} & 当前 hart 的栈或私有堆 & 同一 hart 内；事件链表等待，不使用原子指令。 \\
  \code{dayan_shared_resource_t} & \code{shared} allocation & 跨 hart；使用无锁 AMOCAS/AMOADD 状态字和 scheduler poll。 \\
\end{dayantable}

resource 的函数和所有权规则见第~\ref{sec:kernel-030}~节。

\subsection{C++ 范围}

Device 代码可以使用不依赖 Runtime 支持的类、模板和内联泛型。当前没有 C++ exception、
RTTI、标准 C++ 库、全局构造器或析构器。Dayan OS 不执行 \code{.init_array}；含非平凡
静态初始化的对象不能使用。

\end{document}
